\documentclass[10pt]{article}

\usepackage[utf8]{inputenc}

\usepackage[T1]{fontenc}

\usepackage{amsmath}

\usepackage{amsfonts}

\usepackage{amssymb}

\usepackage[version=4]{mhchem}

\usepackage{stmaryrd}

\usepackage{bbold}

\usepackage{graphicx}

\usepackage[export]{adjustbox}

\graphicspath{ {./images/} }



\begin{document}

такभе $n-с о$ п и то н н ы м и р е пі е н и я м и, к-рые при больших временах $(t \rightarrow \pm \infty)$ можно приближенно зашисать в виде суммы $n$ слагаемых $u_{s}(x, t)$, каждое из к-рых характеризуется своеї скоростью $v_{i}$ и положением центра $x_{0 i}^{ \pm}$. Для $n$-солитонного решения набор скоростей до столкновения $(t \rightarrow-\infty)$ и после столкновения $(t \rightarrow+\infty)$ остается неизменным, возникают только сдвиги центров С. $x_{0}^{+} \neq x_{0}^{-} \cdot$. Найдено много нелинейных эволюционных уравнений c двумя независимыми переменными, к-рые обладают решениями с приведенными выше свойствами. Так, С. нелинейного уравнения ШІёдингера



\[

i \psi_{t}=-\psi_{x x}-|\psi|^{2} \psi, \quad \psi \in \mathbb{C}

\]



однознатно определяется четырьмя параметрами; C. синус Гордона уравнения



\[

\varphi_{t t}-\varphi_{x x}+\sin \varphi=0

\]



определяется двумя параметрами $v, x_{0}$



\[

\varphi_{s}=4 \operatorname{arctg}\left[\exp \pm\left(x-v t-x_{0}\right) / \sqrt{1-v^{2}}\right]

\]



и существует двойной С., к-рый определяется четырьмя параметрами; аналогичная ситуация для уравнения Буссинеска



\[

\varphi_{x x}-\varphi_{t t}+\left(\varphi^{2}\right)_{x x}+\varphi_{x x x x}=0

\]



для уравнения Хироты



$\varphi_{t}+i 3 \alpha|\varphi|^{2} \varphi_{x}+\beta \varphi_{x x}+i \sigma \varphi_{x x x}+\left.\delta \varphi\right|^{2} \varphi=0, \alpha \beta=\sigma \delta$, и т. Д. Существуют физически интересные уравнения и с большим числом независимых переменных, к-рые имеют солитонные решения с приведенными выше свойствами. Напр., С. уравнения Кадомцева - Петвиашвилй



\[

\left(u_{t}+6 u u_{x}+u_{x x x}\right)_{x}=u_{y y},

\]



локализованный по $x$ и $y$, равен



\[

u(x, y, t)=2 \frac{\partial^{2}}{\partial x^{2}} \ln \left(\frac{1}{v^{2}}+\left|x+i v y-3 v^{2} t\right|^{2}\right), \quad v \in \mathbb{R} .

\]



В физич. литературе термин «С.» означает частицеподобное решение нелинейных уравнений классич. теории поля, для к-рого плотности энергии и импульса остаются локализованными в окрестности нек-рой точки пространства в любой момент времени. В ряде случаев локализация может иметь место вблизи замкнутых линий, поверхностей. Такие локализованные решения наз. также к и н к а м и, м о но по л я ми. При отыскании таких решений играют роль топологич. соображения, в частности для ряда моделей удается построить ток $J_{\mu}(x)$, дивергенция к-рого равна нулю независимо от уравнений движения, а соответствующий интеграл движения (топологич. заряд) $Q=\int J_{0}(x) d^{3} x$ оценивает снизу функционал энергии.



\begin{center}

\includegraphics[max width=\textwidth]{2023_07_09_eaff7f880c8955b51f71g-1(2)}

\end{center}



\begin{center}

\includegraphics[max width=\textwidth]{2023_07_09_eaff7f880c8955b51f71g-1}

\end{center}



\includegraphics[max width=\textwidth, center]{2023_07_09_eaff7f880c8955b51f71g-1(5)}

Ду б р о в и н Б. А., М а т в е е в В. Б., Н 0 в

«спехи матем. наук», 1976, т. 31 , в. 1, с. $55-136$.



СОНИНА ИІТЕГРАЛ - представление цилиндрич. функции интегралом по контуру



\[

J_{v}(z)=\frac{1}{2 \pi i} \int_{-\infty}^{0} e^{\frac{z}{2}\left(t-\frac{1}{t}\right) t^{-v-1}} d t,

\]



где $v-$ произвольно, $\operatorname{Re} z>0$ или $-\frac{\pi}{2}<\arg z<\frac{\pi}{2} \quad$ Интеграл этого типа рассмотрен Н. Я. Сониным (1870).



Иногда С. и. называгот интегралі вида:



$J_{m+n+1}(x)=\frac{x^{n+1}}{2^{n} \Gamma(n+1)} \int_{0}^{\infty} J_{m}(x \sin t) \sin ^{m+1} t \cos ^{2 n+1} t d t$ $m, n>-1$. Jlum.: [1] Л а в р е н т ь е в M. А., III а б а т Б. B., Методы теории функций комплексного переменного, 4изд., М., 1973 ; Формулы, графики, таблицы, 2 изд., пер. с нем., М., 1968.



А. Б. Иванов. СООБЩЕНИЙ КВАНТОВАНИЕ - разбиение множества возможных сообщений, вырабатываемых источником сообщений, на конечное (иногда счетное) число неперекрывающихся подмножеств $A_{i}$ так, чтобы сообщения каждого класса могли быть представлены с заданної сообщений точностью воспроизведения нек-рым специально выбранным элементом этого подмножества $a_{i} \in A_{i}$. Заданному С. к. соответствует способ кодйрования источника сообщений, задаваемый кодирующей функцией $\varphi(x)=a_{i}$, если $x \in A_{i}$. Квантование позволяет заменить передачу непрерывного сигнала дискретным сигналом без нарушения условий точности воспроизведения сообщений.



Jum.: [1] Х а р к е в ич А. А., Борьба с помехами, 2 изд., М., 1965; [2] ШІ е н н о н К., Работы по теории информации и кибернетике, пер. с англ., М., 1963; [3] Г а л л а г е p Р., Тео-

рия информации и надежная связь, пер. с англ., М., 1974; [4] B e r g e r T., Rate distortion theory, N. Y., 1971 .



СООБШЕНИИ СКОРОСТЬ СОЗДАНИЯ - величина, характеризующая информации количество, создаваемое за единицу времени источником сообщений. С. с. с. источника сообщений $u$ с дискретным временем, вырабатывающего сообщение $\xi=\left(\ldots, \xi_{-1}, \xi_{0}, \xi_{1}, \ldots\right)$, образованное последовательностью случайных величин $\left\{\xi_{k}, k=\ldots,-1,0,1, \ldots\right\}$, принимающих значение из нек-рого дискретного множества $X$, определяется равенством



\[

\bar{H}(u)=\lim _{n-k \rightarrow \infty} \frac{1}{n-\bar{k}} H\left(\xi_{k}^{n}\right),

\]



если такой предел существует. Здесь $H\left(\xi_{k}^{n}\right)$ - энтропия случайной величины $\xi_{k}^{n}=\left(\xi_{k}, \ldots, \xi_{n}\right)$. Величину $\overline{I I}(u)$, определяемую равенством (*), называют также әнтропией (на символ) источника сообщений.



Доказать существование предела (*) и явно его вычислить удается, напр., для стационарных источников; явные формулы для $\bar{H}(u)$ получены для стационарных марковских источников и гауссовских источников. Понятие С. с. с. $\bar{H}(u)$ тесно связано с понятием избыточности источника сообений.



Если $u$ - стационарный эргодич. истотник сообцений $c$ конечным числом состояний, то справедливо следующее с в о й с т в о а с и м п т о т и ч е с к о ӥ



\includegraphics[max width=\textwidth, center]{2023_07_09_eaff7f880c8955b51f71g-1(4)}

М а кми л л а н a, [1]). Пусть $P\left(x^{L}\right)=\mathrm{P}\left\{\xi^{L}=x^{L}\right\}$, где $x^{L}=\left(x_{1}, \ldots, x_{L}\right)$ суть значения $\xi^{L}=\left(\xi_{1}, \ldots, \xi_{I}\right) \cdots$ отрезка сообщениї длины $L$. Для произвольных $\varepsilon>0$, $\delta>0$ существует $L_{0}(\varepsilon, \delta)$ такое, что при всех $L \geqslant L_{0}(\varepsilon, \delta)$



\[

\mathrm{P}\left\{\left|\frac{-\log P\left(\xi^{L}\right)}{L}-\bar{H}(u)\right|>\delta\right\}<\varepsilon .

\]



Лuт.: [1] В о л ь ф о в и ц Д ж., Теоремы кодирования тео-



\includegraphics[max width=\textwidth, center]{2023_07_09_eaff7f880c8955b51f71g-1(1)}

Теория информации и надежная связь, пер. с англ., М., 1974;



\includegraphics[max width=\textwidth, center]{2023_07_09_eaff7f880c8955b51f71g-1(3)}

СООБШЕНИЙ ТОЧНОСТЬ ВОСПРОИЗВЕДЕНИЯ мера качества передачи сообщений от источника сообцений к получателю (адресату) по каналу связи. Требования, предъявляемые к С. т. в. в теории передачи информации, обычно трактуют статистически, выделяя класс $W$ допустимых совместных распределений для пары $(\xi, \tilde{\xi})$ в множестве всех вероятностных мер в произведении $\left(\mathfrak{X} \times \tilde{\mathfrak{X}}, \quad S_{\mathfrak{X}} \times S_{\mathfrak{X}}\right)$, где $\left(\mathfrak{X}, \quad S_{\mathfrak{X}}\right)$ - измеримое пространство значений сообщения $\xi$, вырабатываемого источником, а $\left(\tilde{\mathfrak{F}}, S_{\tilde{\mathfrak{X}}}\right)$ - измеримое пространство значений воспроизводимого сообщения $\tilde{\xi}$, получасмого адресатом. С. т. в. тасто задают шри помощи меры иска-





\end{document}